\documentclass[11pt]{article}
\usepackage{amsmath,amssymb,amsthm,booktabs,array,geometry,hyperref}
\geometry{margin=1in}

\newtheorem{theorem}{Theorem}
\newtheorem{lemma}[theorem]{Lemma}
\newtheorem{proposition}[theorem]{Proposition}
\newtheorem{corollary}[theorem]{Corollary}
\newtheorem{definition}[theorem]{Definition}
\newtheorem{remark}[theorem]{Remark}
\newtheorem{conjecture}[theorem]{Conjecture}

\title{CONJ\_BOUNDARY\_DECIDABLE and the Schanuel Connection:\\
Logical Dependence of ELC Membership Decidability}
\author{Monogate Research}
\date{2026-04-21}

\begin{document}
\maketitle

\begin{abstract}
We explore CONJ\_BOUNDARY\_DECIDABLE: does deciding whether a given elementary expression
$E$ lies in the real ELC field $\varepsilon(\mathbb{R})$ reduce to Schanuel's conjecture?
We map the exact logical structure: (1) Richardson's undecidability theorem for
expression zero-testing limits the naive approach; (2) conditional on Schanuel,
zero-testing of ELC expressions \emph{is} decidable (Baker-Risch algorithm sketch);
(3) ELC membership is equivalent to a specific zero-testing problem;
(4) therefore, Schanuel implies CONJ\_BOUNDARY\_DECIDABLE for a natural class of inputs;
(5) the converse direction (CONJ\_BOUNDARY\_DECIDABLE implies progress on SC) is also explored.
We conclude with clarifying lemmas and the precise boundary of what SC settles.
\end{abstract}

\section{Problem Formulation}

\begin{definition}[ELC membership problem]
Given a closed-form expression $E$ built from rational constants, variables, $+$, $-$,
$\times$, $\div$, $\exp$, $\ln$, $\sin$, $\cos$, $\int$ (elementary integrals),
and evaluated at a rational point $x_0\in\mathbb{Q}$, decide:
\[
\mathrm{ELC-MEMBER}(E,x_0): \quad E(x_0)\in\varepsilon(\mathbb{R})?
\]
\end{definition}

\begin{remark}
When $E$ uses only $\exp$ and $\ln$ (no trig, no integrals), $E(x_0)\in\varepsilon(\mathbb{R})$
trivially by the closure of $\varepsilon(\mathbb{R})$ under exp and ln. The problem is nontrivial
only when $E$ involves operations outside $\varepsilon(\mathbb{R})$: trig functions, $\int$,
$\Gamma$, etc.
\end{remark}

\section{Richardson's Undecidability and Its Scope}

\begin{theorem}[Richardson, 1968]
\label{thm:richardson}
The problem of deciding whether an expression $E$ built from
$\{+, -, \times, \exp, \sin, |x|, \text{integer constants}\}$
satisfies $E(x)=0$ for all $x$ is undecidable.
\end{theorem}

\begin{remark}
Richardson's result uses $\sin$ and $|x|$ (the absolute value) in an essential way.
If either $\sin$ or $|x|$ is removed from the grammar, the zero-equivalence problem
becomes decidable (this follows from the Risch algorithm for exp/log expressions,
and from the Tarski--Seidenberg theorem for polynomial and semi-algebraic expressions).
\end{remark}

\begin{proposition}[Scope of Richardson for ELC membership]
The ELC membership problem is \emph{not} directly blocked by Richardson's theorem,
because:
\begin{enumerate}
  \item ELC membership asks about a specific value $E(x_0)$, not universal equivalence.
  \item The expressions we are testing (erf, $\ln\Gamma$, $\mathrm{Li}_2$, $\sin(1)$)
    are specific numbers, not universally quantified expressions.
  \item Point-evaluation decidability is weaker than global zero-equivalence.
\end{enumerate}
\end{proposition}

\section{Zero-Testing ELC Expressions Under Schanuel}

\subsection{The Shapiro--Risch--Baker Framework}

\begin{definition}[ELC expression at a rational point]
An \emph{ELC number} is a value of the form $f(x_0)$ where $f\in\mathrm{ELC}(\mathbb{R})$
and $x_0\in\mathbb{Q}$. The set of all ELC numbers is $\varepsilon(\mathbb{R})\cap\mathbb{R}$
(a subfield of $\mathbb{R}$).
\end{definition}

\begin{theorem}[Zero-testing under SC --- Baker-type result]
\label{thm:zeroing}
Assume Schanuel's conjecture. Then:
\begin{enumerate}
  \item Any nonzero ELC number $c\in\varepsilon(\mathbb{R})$ satisfies $|c|>\exp(-H(c)^C)$
    for some effective constants depending on the expression complexity $H(c)$
    (an effective ``lower bound'' analogous to Baker's theorem for linear forms in logarithms).
  \item Given an ELC expression $E$ with rational constants, one can algorithmically
    decide whether $E(x_0)=0$ for $x_0\in\mathbb{Q}$: enumerate the ELC expression tree
    and apply the lower bound to eliminate zero vs.\ nonzero.
\end{enumerate}
\end{theorem}

\begin{proof}[Sketch]
Under SC, the transcendence degree of any finite set of exponentials and logarithms
over $\mathbb{Q}$ is at least the linear independence dimension of the arguments.
This gives effective algebraic independence measures, which translate to effective
lower bounds on nonzero expressions via a compactness argument (the expression either
cancels algebraically---which is checkable---or is bounded away from $0$ by the SC measure).
\end{proof}

\section{ELC Membership as a Zero-Testing Problem}

\begin{theorem}[Reduction of ELC membership to zero-testing]
\label{thm:reduction}
$E(x_0)\in\varepsilon(\mathbb{R})$ if and only if there exists an ELC expression $F$
such that $E(x_0)-F(x_0)=0$.
\end{theorem}

\begin{proof}
Trivial in one direction. In the other: if $E(x_0)=c\in\varepsilon(\mathbb{R})$,
then there exists an ELC expression $F$ (possibly the constant $c$ expressed as an ELC tree)
with $F(x_0)=c$, giving $E(x_0)-F(x_0)=0$.
\end{proof}

\begin{remark}
The non-trivial content of CONJ\_BOUNDARY\_DECIDABLE is that there is an \emph{effective}
algorithm to search for or rule out such $F$. The difficulty:
\begin{enumerate}
  \item The space of ELC expressions $F$ is infinite (all finite-depth ELC trees).
  \item We need to either find $F$ or prove no such $F$ exists.
\end{enumerate}
Under Theorem~\ref{thm:zeroing} (SC), we can in principle test:
for each candidate $F$ (enumerated by depth), decide whether $E(x_0)-F(x_0)=0$
using the SC-based zero-test. If the answer is always ``No'' and we can bound
the depth at which a match must occur (if it occurs), we obtain decidability.
\end{remark}

\section{The Schanuel Implication: SC $\Rightarrow$ CONJ\_BOUNDARY\_DECIDABLE}

\begin{theorem}[SC implies conditional decidability]
\label{thm:sctocbd}
Assume:
\begin{enumerate}
  \item[(A)] Schanuel's conjecture (SC).
  \item[(B)] The expression $E$ is built from the 23-operator set plus trig functions $\sin,\cos$
    and a finite list of transcendental integrals (erf, $\ln\Gamma$, $\mathrm{Li}_2$, etc.).
  \item[(C)] $x_0\in\mathbb{Q}$.
\end{enumerate}
Then ELC-MEMBER$(E,x_0)$ is decidable.
\end{theorem}

\begin{proof}[Sketch]
Under SC + (B), every nonzero combination $E(x_0)-F(x_0)$ for $F\in\mathrm{ELC}$
satisfies an effective lower bound $|E(x_0)-F(x_0)|\geq \delta(H)>0$, where $H$
is the complexity (number of nodes, height of coefficients) of the combined expression.

Algorithm:
\begin{enumerate}
  \item Enumerate ELC expressions $F_1,F_2,\ldots$ in order of increasing depth.
  \item For each $F_k$: compute $d_k=E(x_0)-F_k(x_0)$ to precision $\delta(H_k)/2$.
  \item If $|d_k|<\delta(H_k)/2$: by the SC lower bound, $d_k=0$, so $E(x_0)=F_k(x_0)\in\varepsilon(\mathbb{R})$. Return YES.
  \item Otherwise: $d_k\neq0$; continue.
  \item Termination: if $E(x_0)\notin\varepsilon(\mathbb{R})$, then by the AIL/SC framework,
    no $F_k$ will match. We need a termination criterion for the NO case.
\end{enumerate}

The termination issue: without an upper bound on the depth of the matching $F$,
the algorithm is semi-decidable (YES is always found in finite time; NO may not terminate).
Under SC + a strengthening (an \emph{effective depth bound}: if $E(x_0)\in\varepsilon(\mathbb{R})$
with an $n$-node ELC expression, then there is an expression of depth $\leq B(H_E)$ for
some effective bound $B$), the algorithm terminates.
\end{proof}

\begin{remark}
The \emph{effective depth bound} is the crucial missing piece.
It is not implied by SC alone; it requires an additional ``complexity'' version of SC
that bounds the minimum ELC tree representing a given ELC number in terms of its
arithmetic complexity. This is analogous to the Kolmogorov complexity interpretation
of SC and is an open problem.
\end{remark}

\section{The Converse: CONJ\_BOUNDARY\_DECIDABLE $\Rightarrow$ ? }

\begin{proposition}[What decidability implies]
If ELC-MEMBER is decidable (CONJ\_BOUNDARY\_DECIDABLE), then:
\begin{enumerate}
  \item For every algebraic $\alpha\in\overline{\mathbb{Q}}$, $\sin(\alpha)\notin\varepsilon(\mathbb{R})$
    can be decided: this is already handled by HLW (unconditional).
  \item For every specific transcendental constant ($\pi$, $e^\pi$, $\Gamma(1/4)$, etc.),
    membership in $\varepsilon(\mathbb{R})$ is decidable.
  \item The ``constant problem''---deciding whether two ELC numbers are equal---is decidable
    (by ELC-MEMBER applied to their difference).
\end{enumerate}
\end{proposition}

\begin{theorem}[Partial SC-consequence of decidability]
\label{thm:converse}
If ELC-MEMBER is decidable and the algorithm is \emph{uniform} in the following sense:
for each expression $E$ it produces either a matching ELC expression $F$ or a proof
that none exists (not just an answer YES/NO), then the algorithm witnesses effective
lower bounds on nonzero ELC differences.
Such effective lower bounds, when extracted and generalized, imply specific cases of SC
(namely, for expressions in the decidability grammar).
\end{theorem}

\begin{proof}
The ``proof of NO'' in a decidability algorithm must certify $E(x_0)\notin\varepsilon(\mathbb{R})$.
By Theorem~\ref{thm:reduction}, this is equivalent to certifying $E(x_0)-F(x_0)\neq0$
for all ELC $F$. For the certificate to be finite and checkable, it must invoke
an algebraic independence result (an explicit lower bound). Extracting this bound
gives a specific case of SC: the transcendence degree of the expressions involved
is at least the linear independence dimension of their arguments.
\end{proof}

\section{Cases Where Schanuel Is Not Needed}

\begin{proposition}[Decidable without SC]
The following cases of ELC-MEMBER are decidable unconditionally:
\begin{enumerate}
  \item $E$ uses only $\exp$ and $\ln$: trivially in $\varepsilon(\mathbb{R})$. Decide YES.
  \item $E=\sin(r)$ for $r\in\mathbb{Q}$: $\sin(r)\notin\varepsilon(\mathbb{R})$ iff $r\neq0$
    (by HLW). Decidable: YES iff $r=0$.
  \item $E=p(\sin(r_1),\ldots,\sin(r_k))$ for $p$ a polynomial with rational coefficients
    and $r_i\in\mathbb{Q}$: decidable by the algebraic independence of $\sin(r_i)$ over $\mathbb{Q}$
    (HLW implies these are independent transcendentals when the $r_i/\pi$ are linearly
    independent over $\mathbb{Q}$).
  \item $E=\mathrm{erf}(1)$: Risch algorithm certifies $\notin$ elementary functions,
    hence $\notin\varepsilon(\mathbb{R})$. Decide NO.
  \item $E=\ln\Gamma(p/q)$ for $p/q\notin\mathbb{Z}^+$: Gauss's reflection formula and
    the algebraic independence of $\pi$ certify $\notin\varepsilon(\mathbb{R})$. Decide NO.
\end{enumerate}
\end{proposition}

\section{The Exact Logical Dependence Map}

\begin{center}
\renewcommand{\arraystretch}{1.3}
\begin{tabular}{lll}
\toprule
Statement & Strength & Status \\
\midrule
HLW (Hermite--Lindemann--Weierstrass) & Decidability for algebraic-argument trig & Proved \\
Risch algorithm & Decidability for elementary integrals & Proved \\
Baker's theorem & Effective lower bounds for log-linear forms & Proved \\
Gelfond--Schneider & Transcendence of $\alpha^\beta$, $\alpha,\beta\in\overline{\mathbb{Q}}$ & Proved \\
\midrule
AIL (our algebraic-independence lemma) & $\sin(1)\notin\varepsilon(\mathbb{R})$ & Proved \\
\midrule
Schanuel's conjecture (SC) & General transcendence degree lower bound & Open \\
GAIL (generalized AIL) & $\sin(\alpha)\notin\varepsilon(\mathbb{R})$ for $\alpha\notin\varepsilon(\mathbb{R})$ & Open; follows from SC \\
Effective depth bound (EDB) & Matching ELC $F$ has bounded depth & Open; strengthens SC \\
\midrule
CONJ\_BOUNDARY\_DECIDABLE & ELC-MEMBER is decidable & Open; follows from SC+EDB \\
Zero-testing of ELC expressions & Equality of two ELC numbers & Open; follows from SC \\
\bottomrule
\end{tabular}
\end{center}

\begin{theorem}[Main reduction]
\label{thm:main}
The following implications hold:
\[
\mathrm{SC} + \mathrm{EDB} \Rightarrow \mathrm{CONJ\_BOUNDARY\_DECIDABLE} \Rightarrow \text{specific cases of SC}.
\]
In particular:
\begin{itemize}
  \item SC alone implies semi-decidability (YES is always verified in finite time).
  \item SC + EDB implies full decidability.
  \item CONJ\_BOUNDARY\_DECIDABLE (as a uniform algorithm) implies specific instances of SC.
  \item CONJ\_BOUNDARY\_DECIDABLE does \emph{not} imply SC in general without additional
    uniformity assumptions.
\end{itemize}
\end{theorem}

\section{Clarifying Lemmas}

\begin{lemma}[ELC field is not effectively enumerable in the usual sense]
$\varepsilon(\mathbb{R})$ is a countable subfield of $\mathbb{R}$, but the \emph{membership test}
(given a real number $r$, is $r\in\varepsilon(\mathbb{R})$?) is not decidable from a raw
decimal expansion: we cannot in general recognize an ELC number from its decimal expansion alone.
\end{lemma}

\begin{proof}
Any effective recognition of ELC numbers from decimal expansions would give a decision
procedure for the zero-equivalence problem (compare two ELC number decimal expansions),
which is undecidable by Richardson's theorem for sufficiently rich expression grammars.
However, our ELC grammar (exp + log, no $\sin$) has decidable zero-equivalence (Risch),
so the lemma applies only to grammars including trig or transcendental integrals.
\end{proof}

\begin{lemma}[SC is not equivalent to CONJ\_BOUNDARY\_DECIDABLE]
SC is strictly stronger than CONJ\_BOUNDARY\_DECIDABLE: there exist statements implied
by SC that are not implied by ELC-MEMBER decidability alone.
\end{lemma}

\begin{proof}
SC implies algebraic independence of $\{\pi, e^\pi, e^{e^\pi}, \ldots\}$ (Nesterenko's theorem
implies some of this). CONJ\_BOUNDARY\_DECIDABLE decides only whether a specific value
is in $\varepsilon(\mathbb{R})$; it does not give the full transcendence degree of
arbitrary exponential-logarithmic expressions.
\end{proof}

\begin{lemma}[The restriction to real ELC is essential]
CONJ\_BOUNDARY\_DECIDABLE for real $\varepsilon(\mathbb{R})$ is strictly easier than
the analogous problem for complex $\varepsilon(\mathbb{C})$, because:
\begin{enumerate}
  \item The complex ELC field includes $\sin(1)=\mathrm{Im}(e^i)$ as a 1-node expression.
  \item Deciding real vs.\ complex ELC membership for trig values is essentially the
    imaginary-part projection problem.
  \item SC for real ELC restricts to real arguments and real exponentials, a proper subcase.
\end{enumerate}
\end{lemma}

\section{Summary and Open Questions}

\begin{enumerate}
  \item \textbf{SC $\Rightarrow$ semi-decidability}: Under SC, ELC-MEMBER is semi-decidable (YES in finite time).
  \item \textbf{SC + EDB $\Rightarrow$ full decidability}: With an effective depth bound, ELC-MEMBER is fully decidable.
  \item \textbf{Converse (partial)}: A uniform YES/NO algorithm for ELC-MEMBER implies specific instances of SC.
  \item \textbf{Richardson limit}: Full decidability for grammars including $\sin$ and $|x|$ is blocked by Richardson; the ELC membership problem avoids this by restricting to specific function classes.
  \item \textbf{Already decidable}: All cases with algebraic-argument trig (HLW), elementary integrals (Risch), and specific constants ($\pi$, $e^\pi$) are decidable unconditionally.
  \item \textbf{Open}: The effective depth bound (EDB). Without it, SC gives only semi-decidability.
  \item \textbf{Open}: Whether CONJ\_BOUNDARY\_DECIDABLE follows from SC alone (without EDB).
\end{enumerate}

\begin{conjecture}[Refined CONJ\_BOUNDARY\_DECIDABLE]
ELC-MEMBER is decidable for all inputs in the grammar of Section 1, assuming:
(a) Schanuel's conjecture, and
(b) the effective depth bound: if $E(x_0)\in\varepsilon(\mathbb{R})$, then the minimum ELC tree
for $E(x_0)$ has depth at most $B(H_E)$ for a primitive recursive function $B$.
Part (b) is expected to hold and may follow from SC + an effective version of
the Risch structure theorem, but is not yet proved.
\end{conjecture}

\end{document}
